% (The MIT License)
%
% Copyright (c) 2021 Yegor Bugayenko
%
% Permission is hereby granted, free of charge, to any person obtaining a copy
% of this software and associated documentation files (the 'Software'), to deal
% in the Software without restriction, including without limitation the rights
% to use, copy, modify, merge, publish, distribute, sublicense, and/or sell
% copies of the Software, and to permit persons to whom the Software is
% furnished to do so, subject to the following conditions:
%
% The above copyright notice and this permission notice shall be included in all
% copies or substantial portions of the Software.
%
% THE SOFTWARE IS PROVIDED 'AS IS', WITHOUT WARRANTY OF ANY KIND, EXPRESS OR
% IMPLIED, INCLUDING BUT NOT LIMITED TO THE WARRANTIES OF MERCHANTABILITY,
% FITNESS FOR A PARTICULAR PURPOSE AND NONINFRINGEMENT. IN NO EVENT SHALL THE
% AUTHORS OR COPYRIGHT HOLDERS BE LIABLE FOR ANY CLAIM, DAMAGES OR OTHER
% LIABILITY, WHETHER IN AN ACTION OF CONTRACT, TORT OR OTHERWISE, ARISING FROM,
% OUT OF OR IN CONNECTION WITH THE SOFTWARE OR THE USE OR OTHER DEALINGS IN THE
% SOFTWARE.

\documentclass[compact,manuscript]{./ybook}
\renewcommand*\theversion{0.0.0}
\renewcommand*\thedate{00.00.0000}
\renewcommand*\thetitle{\LaTeX{} Class \ff{ybook}}
\renewcommand*\theauthor{\nospell{Yegor Bugayenko}}
\begin{document}

\includegraphics[height=6em]{images/signature-black.pdf}
\vspace*{24pt}

\maketitle

\index{Amazon}
\first{T}{he provided} class \ff{ybook} helps me design
\href{https://www.yegor256.com/books.html}{my books} and
publish them
\href{https://www.amazon.com/Yegor-Bugayenko/e/B01AM1QMDK}{on Amazon}.
You are welcome to use it too, it's pretty easy:

\begin{listing}
\documentclass[draft]{ybook}
\begin{document}
Hello, world!
\end{document}
\end{listing}

There are a few class options you can use:

\ff{compact} --- when you need to make text more compact
and take less vertical space;

\ff{manuscript} --- when the format is not for Amazon printing,
but for some other purposes (the page size is A4), like the document
you are reading right now;

\ff{draft} --- when it's a draft for reviewers (the page size is A4)
and you want to have a watermark and a compact form of the content;

\ff{cover=book.pdf} --- when you print a cover and the book is in the
\ff{book.pdf} file.

\section*{Printers}

There are a number of printers, commands that print large blocks of text
in the expected format:

\ff{\textbackslash{}ybPrintTitlePage} prints the first page of a book.

\ff{\textbackslash{}ybPrintTOC} prints the table of contents.

\ff{\textbackslash{}ybPrintBibliography} prints the list of bib references.

\ff{\textbackslash{}ybPrintIcons} prints your social icons.

\ff{\textbackslash{}ybPrintIndex[1]} prints index with an optional name of the section (instead of ``Index'').

\ff{\textbackslash{}ybPrintCover} prints the cover of the book.

For more samples, check the directory \ff{samples/}.

\section*{Cover Page}

\first{A}{ cover} page must have a special formatting. You do it this way
(provided your book already exists in the \ff{book.pdf} file):

\begin{listing}
\documentclass[cover=book.pdf]{ybook}
\renewcommand*\theauthor{Jeff Lebowski}
\renewcommand*\theversion{0.1}
\renewcommand*\theyear{2021}
\renewcommand*\thevolume{1}
\renewcommand*\theprice{\$14.99}
\renewcommand*\thespinetext{new book, vol.1}
\renewcommand*\thetitle{A New Book}
\renewcommand*\theimage{images/cactus} % file name of PDF image
\renewcommand*\thetldr{Put your TL;DR text here...}
\begin{document}
\ybPrintCover
\end{document}\end{listing}

The PDF file generated most likely will be accepted by Amazon without
any questions. The \ff{\textbackslash{}ybPrintCover} command will
make sure the size of the page is correctly adjusted for the
number of pages in the \ff{book.pdf} document.

\section*{Commands}

\first{T}{here} are a few new commands you may want to learn and use:

\index{font}
\ff{\textbackslash{}ff[1]} makes your code fixed-font, like \ff{this one}.

\ff{\textbackslash{}tbd[1]} makes some text \tbd{highlighted}, in order to indicate
that it will be written later.

\ff{\textbackslash{}first[2]} makes the first letter large and the
rest of the text small caps, in order to emphasize the start of a new
group of paragraphs.

You are welcome to suggest additional commands, but the style
of my books is intentionally as simple as possible, avoiding formatting
as much as possible.

\end{document}